\subsection{Okluzja jako źródło utraty informacji i szumu w rozpoznawaniu twarzy}

\paragraph{}Fundamentem dla zrozumienia metod radzenia sobie z niepełnymi danymi wizualnymi jest zdefiniowanie natury problemu okluzji. Jak wskazują Zeng i in. [1] w swoim przeglądzie technik rozpoznawania twarzy, okluzja stanowi dwojakie wyzwanie dla systemów biometrycznych. Po pierwsze, powoduje utratę informacji dyskryminacyjnych w obszarach zasłoniętych. Po drugie, i co często groźniejsze, wprowadza szum (błędne cechy) pochodzący z obiektu przysłaniającego (np. maski, dłoni), który może zostać błędnie zinterpretowany przez sieć jako cecha tożsamościowa.

Zeng i in. proponują systematykę podejść badawczych, dzieląc je na trzy główne kategorie: (1) odporną ekstrakcję cech (ang. Occlusion Robust Feature Extraction), która dąży do uzyskania reprezentacji niewrażliwej na zakłócenia, często poprzez skupienie się na cechach lokalnych zamiast globalnych; (2) rozpoznawanie świadome okluzji (ang. Occlusion Aware Face Recognition), gdzie mechanizmy detekcji lub maskowania dynamicznie wykluczają zaszumione obszary z procesu decyzyjnego; oraz (3) rekonstrukcję okluzji (ang. Occlusion Recovery), polegającą na wizualnym odtworzeniu brakujących fragmentów [1].

W kontekście tej klasyfikacji, współczesna literatura coraz częściej odchodzi od prostych metod rekonstrukcji na rzecz zaawansowanego modelowania relacji przestrzennych (podejście typu 1 i 2). Klasyczne sieci splotowe (CNN), choć skuteczne w ekstrakcji cech lokalnych, mogą napotykać trudności w modelowaniu globalnych zależności geometrycznych, gdy część danych wejściowych ulega zmianie